\section*{第 \ref{ch:g8:from-egg-to-birth} 章——从卵细胞到出生}

\begin{solution}{exo:g8:from-egg-to-birth:1}
\omterm{def:g5:flower-to-fruit:steps}{受精}后一天内开始分裂；第四天时形成
由数十个\omterm{def:g5:first-look-at-cells:cell}{细胞}组成的桑葚胚，沿\omterm{def:g8:reproductive-systems:female}{输卵管}下行；
约第七天时着床于\omterm{prop:g5:human-reproduction:plan}{子宫}内膜。
\end{solution}

\begin{solution}{exo:g8:from-egg-to-birth:2}
最初依靠\omterm{def:g4:animal-reproduction:sexual}{卵细胞}自身的储备，即大型\omterm{def:g8:sexual-reproduction:gametes}{配子}形成时
所储存的营养。着床结束了这种状况：
\omterm{def:g5:human-reproduction:pregnancy}{胚胎}接入\omterm{prop:g5:human-reproduction:plan}{子宫}内膜的供应，
不久便组织成\omterm{def:g8:from-egg-to-birth:placenta}{胎盘}。
\end{solution}

\begin{solution}{exo:g8:from-egg-to-birth:3}
嵌入\omterm{prop:g5:human-reproduction:plan}{子宫}壁的“生命支持盘”，即\omterm{def:g8:from-egg-to-birth:placenta}{胎盘}，
母体与\omterm{def:g5:human-reproduction:pregnancy}{胚胎}的\omterm{prop:g4:journey-of-food:absorb}{血液}在此靠近却互不混合。
流向\omterm{def:g5:human-reproduction:pregnancy}{胚胎}的是\omterm{def:g7:digestion-and-nutrients:families}{营养物质}和\omterm{def:g7:what-is-respiration:respiration}{氧气}；反方向流出的是废物和\omterm{def:g7:what-is-respiration:respiration}{二氧化碳}。
两种\omterm{prop:g4:journey-of-food:absorb}{血液}本身从不相互穿越。
\end{solution}

\begin{solution}{exo:g8:from-egg-to-birth:4}
第 1--8 周为\omterm{def:g5:human-reproduction:pregnancy}{胚胎}期：由分裂、迁移、分化的\omterm{def:g5:first-look-at-cells:cell}{细胞}
搭建出身体的基本构造。第 9--40 周为\omterm{def:g5:human-reproduction:pregnancy}{胎儿}期：
在此基础上继续生长发育，脑部\omterm{def:g5:brain-in-command:system}{神经}连接逐步形成，
肺部开始“预演”\omterm{def:g7:what-is-respiration:respiration}{呼吸}运动。
\end{solution}

\begin{solution}{exo:g8:from-egg-to-birth:5}
因为整个发育过程依赖物质的输送：\omterm{def:g8:from-egg-to-birth:placenta}{胎盘}的
物质交换一旦开始，就必须有泵推动\omterm{def:g5:human-reproduction:pregnancy}{胚胎}的\omterm{prop:g4:journey-of-food:absorb}{血液}
经\omterm{prop:g5:human-reproduction:cord}{脐带}和身体循环，这条供应线
不能等待发育计划的其余部分。
\end{solution}

\begin{solution}{exo:g8:from-egg-to-birth:6}
分娩时\omterm{prop:g5:human-reproduction:plan}{子宫}收缩逐渐增强；\omterm{def:g5:human-reproduction:pregnancy}{胎儿}通常头先娩出；
转换过程：婴儿第一声啼哭，\omterm{prop:g5:human-reproduction:cord}{脐带}剪断，
新生儿自身的交换界面接管工作；
\omterm{def:g8:from-egg-to-birth:placenta}{胎盘}娩出，即\omterm{def:g8:from-egg-to-birth:placenta}{胎盘}与\omterm{prop:g5:human-reproduction:plan}{子宫}分离并随之排出。
\end{solution}

\begin{solution}{exo:g8:from-egg-to-birth:7}
肺的功能（\omterm{def:g7:what-is-respiration:respiration}{氧气}进入、\omterm{def:g7:what-is-respiration:respiration}{二氧化碳}排出，
参见\cref{ch:g7:lungs-and-gas-exchange}），肠道的功能（\omterm{def:g7:digestion-and-nutrients:families}{营养物质}进入，
参见\cref{ch:g7:digestion-and-nutrients}），
以及\omterm{def:g7:removing-wastes:kidneys}{肾脏}的功能（废物排出，
参见\cref{ch:g7:removing-wastes}）。三个交换界面，共用一个\omterm{def:g8:from-egg-to-birth:placenta}{胎盘}。
\end{solution}

\begin{solution}{exo:g8:from-egg-to-birth:8}
因为搭建计划的几周里，一条受干扰的指令
会在整个发育过程中产生连锁影响——每个器官
都是由一小群正在分化的\omterm{def:g5:first-look-at-cells:cell}{细胞}构成的；又因为这几周
可能在\omterm{def:g5:human-reproduction:pregnancy}{怀孕}尚未被察觉时就已过去——
所以一旦有\omterm{def:g5:human-reproduction:pregnancy}{怀孕}的可能，就应开始护理，
而不是等到第一次超声检查。
\end{solution}

\begin{solution}{exo:g8:from-egg-to-birth:9}
\omterm{def:g7:what-is-respiration:respiration}{呼吸}与啼哭：肺的交换界面，出生仅几分钟。
心跳：驱动两条回路的泵。肤色：输送服务真正
开始输送——富含氧的\omterm{prop:g4:journey-of-food:absorb}{血液}
抵达\omterm{def:g1:the-five-senses:sense}{皮肤}\omterm{def:g7:heart-and-circulation:vessels}{毛细血管}。运动完成了“指挥官”
的检查；评分则是对这一循环的现场审核，
即时进行。
\end{solution}

\begin{solution}{exo:g8:from-egg-to-birth:10}
发生在最初的分裂中——\omterm{prop:g8:sexual-reproduction:core}{受精卵}的球状体
在最初几天内分裂成两个球。两个双胞胎随后
由同一起始\omterm{def:g5:first-look-at-cells:cell}{细胞}所继承的遗传信息发育而来，
（参见\cref{ex:g5:first-look-at-cells:one}），因此两者
遵循相同的发育方案：是复制品，因为分叉发生在
遗传信息确定之后。
\end{solution}

\begin{solution}{exo:g8:from-egg-to-birth:11}
\omterm{def:g5:first-look-at-cells:cell}{细胞}：\omterm{def:g5:human-reproduction:pregnancy}{胎儿}期——器官已经形成，正在生长；
\omterm{def:g3:bones-and-muscles:muscle}{肌肉}\omterm{def:g5:first-look-at-cells:cell}{细胞}与\omterm{def:g5:brain-in-command:system}{神经}\omterm{def:g5:first-look-at-cells:cell}{细胞}正在相互连接。
供应：\omterm{def:g8:from-egg-to-birth:placenta}{胎盘}的物质交换全速运转——\omterm{def:g7:digestion-and-nutrients:families}{营养物质}和\omterm{def:g7:what-is-respiration:respiration}{氧气}
进入，废物排出，母体\omterm{prop:g4:journey-of-food:absorb}{血液}中不允许
有有害物质通过。\omterm{def:g8:puberty:hormones}{激素}：使\omterm{prop:g5:human-reproduction:plan}{子宫}保持安静，
\omterm{prop:g8:reproductive-systems:cycle}{月经周期}暂停。护理：母亲的饮食与休息
要支撑两个交换界面的运转负荷；胎动本身
就是循环回路正常工作的令人安心的信号。
\end{solution}

\begin{solution}{exo:g8:from-egg-to-birth:12}
\omterm{def:g8:from-egg-to-birth:placenta}{胎盘}承担着\omterm{def:g7:what-is-respiration:respiration}{氧气}的摄入、\omterm{def:g7:what-is-respiration:respiration}{二氧化碳}的清除、
\omterm{def:g7:digestion-and-nutrients:families}{营养物质}的供应和废物的清除。第一声啼哭时，
肺接过两种气体的交换；肠道（在首次哺乳后）
接过营养供应；\omterm{def:g7:removing-wastes:kidneys}{肾脏}接过废物清除——
这一切都在\omterm{prop:g5:human-reproduction:cord}{脐带}剪断后几分钟内完成，是一次
界面交接，每个器官都必须一次成功，
这正是产房检查表
所关注的内容。
\end{solution}

\begin{solution}{pb:g8:from-egg-to-birth:1}
\textbf{1.} 着床后开始分泌的妊娠\omterm{def:g8:puberty:hormones}{激素}，
即已安居的\omterm{def:g5:human-reproduction:pregnancy}{胚胎}发出的化学信号；
约从第七天起，就可以检测到它们。

\textbf{2.} \omterm{def:g5:human-reproduction:pregnancy}{胚胎}期结束：第 1--8 周已完成全部
构造计划，每个器官都已开始形成，大小如豆粒。
\omterm{def:g5:human-reproduction:pregnancy}{胎儿}期——即构建与成熟阶段——开始。

\textbf{3.} 酒精和烟雾会穿过\omterm{def:g8:from-egg-to-birth:placenta}{胎盘}进入\omterm{def:g5:human-reproduction:pregnancy}{胎儿}体内，而搭建计划的
几周最容易受到指令干扰；叶酸则是最早的构造过程所需的物质，
而这几周往往在不知不觉中过去，因此
维生素所说的“从开始之前就服用”，
与戒酒、戒烟的道理是同一笔账——
都要在\omterm{def:g5:human-reproduction:pregnancy}{怀孕}被察觉之前就行动。

\textbf{4.} \omterm{def:g8:from-egg-to-birth:placenta}{胎盘}这一“经济体”一开始就需要自己的“泵”：
没有搏动的\omterm{def:g4:heart-and-blood:heart}{心脏}，物资便无法循环，所以它是最先启动的
机器——第 12 周超声影像中的闪动，正是它
运行了八周时的
模样。

\textbf{5.} \omterm{def:g5:human-reproduction:pregnancy}{胎儿}期：已定计划的生长与成熟阶段。\omterm{def:g8:from-egg-to-birth:placenta}{胎盘}被当作
肺（气体交换）、肠（营养供给）和肾（废物清除）来检视——
\omterm{def:g5:human-reproduction:pregnancy}{胎儿}的三道边界，都集中在同一个圆盘上，
此外还要看\omterm{def:g8:from-egg-to-birth:placenta}{胎盘}是否远离出口，
位置是否离出口足够远。

\textbf{6.} 报告—决定—下令的环路：\omterm{def:g3:bones-and-muscles:muscle}{肌肉}
依\omterm{def:g5:brain-in-command:system}{神经}指令收缩——指挥中枢的线路及其“马达”，
借\omterm{prop:g5:human-reproduction:plan}{子宫}壁
显露出来。

\textbf{7.} 因为它同时服务两方的边界：母体自身的，
以及经由\omterm{def:g8:from-egg-to-birth:placenta}{胎盘}的物质交换——\omterm{def:g5:human-reproduction:pregnancy}{胎儿}的。血糖与压强
反映的是双倍负荷
如何被承担。

\textbf{8.} 没有空气的\omterm{def:g7:what-is-respiration:respiration}{呼吸}运动——
风箱空转练习，\omterm{def:g3:bones-and-muscles:muscle}{肌肉}与反射准备
迎接第一声啼哭这一瞬间，那时器官
必须一次便能正常工作。

\textbf{9.} \omterm{prop:g5:human-reproduction:plan}{子宫}——人体最强的\omterm{def:g3:bones-and-muscles:muscle}{肌肉}；
它的任务：打开通道，并将婴儿推向
产道；指令以\omterm{def:g8:puberty:hormones}{激素}形式传递，是婴儿与母亲的
化学物质在相互“对话”。

\textbf{10.} 啼哭：肺已充盈并开始工作——
一道边界已经交接。\omterm{prop:g5:human-reproduction:cord}{脐带}被夹住并剪断：
\omterm{def:g8:from-egg-to-birth:placenta}{胎盘}的物质交换结束。\omterm{def:g8:from-egg-to-birth:placenta}{胎盘}完整娩出：
这一次性使用的器官已完全分离——交接
完成，不留任何残余。

\textbf{11.} 供应与排出的环路——肺、
泵、物资、指挥中枢——独自运行：
本书用两个年级描绘的这台机器，此处首次作一次“审计”。

\textbf{12.} “四十周内，一个\omterm{def:g5:flower-to-fruit:steps}{受精}\omterm{def:g4:animal-reproduction:sexual}{卵细胞}通过分裂、分化
与\omterm{def:g8:puberty:hormones}{激素}调控下的管理，
由一个了不起的圆盘供应，
共建成三万亿个分化\omterm{def:g5:first-look-at-cells:cell}{细胞}。”
\end{solution}
